
\chapter{Modelo de Base de Datos (ERD)}\label{ch:modeloerd}

El diagrama Entidad-Relación (ERD) representa el modelo de datos relacional de FisioManager. A diferencia del diagrama de clases conceptuales (Capítulo~\ref{ch:diagramaclases}), que modela el dominio del problema, el ERD define la estructura física de persistencia: tablas, columnas, tipos de datos, claves primarias, claves foráneas y restricciones de integridad.

\section{Tablas del modelo}

El modelo consta de \textbf{30 tablas} gestionadas con Prisma ORM sobre PostgreSQL, organizadas en siete dominios funcionales. Las claves primarias son identificadores de texto tipo \codeinline{cuid} y los campos con valores acotados se almacenan como \codeinline{TEXT} validados por la aplicación (Zod), no como \codeinline{ENUM} de base de datos.

\subsection{Dominio de Usuarios}

\begin{itemize}
  \item \codeinline{usuarios}: tabla central con los datos comunes a todos los roles (nombre, email único, contraseña hasheada, rol, avatar, teléfono, verificación de email y consentimiento de audio). El campo \codeinline{rol} toma los valores \codeinline{PATIENT}, \codeinline{THERAPIST} y \codeinline{ADMIN}.

  \item \codeinline{perfiles\_paciente}: extensión 1:1 de \codeinline{usuarios} con edad, condición clínica y el terapeuta asignado (\codeinline{terapeuta\_id}). Su PK es \codeinline{usuario\_id}.

  \item \codeinline{perfiles\_terapeuta}: extensión 1:1 de \codeinline{usuarios} con cédula y especialidad.

  \item \codeinline{disponibilidad\_terapeuta}: franjas horarias semanales de atención del terapeuta (día, hora inicio/fin), únicas por día.
\end{itemize}

\subsection{Dominio Clínico}

\begin{itemize}
  \item \codeinline{historial\_clinico}: expediente único por paciente con tipo de sangre, alergias y antecedentes.

  \item \codeinline{diagnosticos}: diagnósticos del historial con código CIE-10, descripción y estado (\codeinline{ACTIVE}/\codeinline{RESOLVED}).

  \item \codeinline{notas\_clinicas}: notas de seguimiento con autor, contenido y bandera de visibilidad para el paciente.

  \item \codeinline{documentos\_clinicos}: archivos adjuntos al expediente (radiografías, laboratorios, consentimientos) con tipo MIME, tamaño, categoría y visibilidad.

  \item \codeinline{planes\_tratamiento}: planes que vinculan paciente y terapeuta, con objetivo clínico, fechas y estado (\codeinline{ACTIVE}, \codeinline{COMPLETED}, \codeinline{SUSPENDED}).

  \item \codeinline{fases\_tratamiento}: fases secuenciales del plan con nombre, orden, duración en semanas y objetivos.
\end{itemize}

\subsection{Dominio de Rutinas y Ejercicios}

\begin{itemize}
  \item \codeinline{plantillas\_actividad}: catálogo de ejercicios reutilizables con tipo (\codeinline{PHYSICAL}/\codeinline{BREATHING}), parte del cuerpo y URLs de video/imagen.

  \item \codeinline{rutinas}: rutinas de ejercicio. \codeinline{paciente\_id} es nullable: \codeinline{NULL} indica plantilla de biblioteca; con valor, es una rutina asignada. Tipo \codeinline{TREATMENT} o \codeinline{RELAXATION}.

  \item \codeinline{actividades}: ejercicios concretos de una rutina (instancian opcionalmente una plantilla) con duración, repeticiones, descanso y orden.

  \item \codeinline{asignaciones\_rutina}: asignación de una rutina a un paciente por un terapeuta, con frecuencia (\codeinline{DAILY}, \codeinline{EVERY\_OTHER\_DAY}, \codeinline{WEEKLY}), fechas, estado y vínculo opcional a una fase.

  \item \codeinline{sesiones} y \codeinline{sesion\_ejercicios}: registro de cada ejecución de una rutina y el detalle por ejercicio (\codeinline{COMPLETED}/\codeinline{SKIPPED}).
\end{itemize}

\subsection{Dominio de Seguimiento}

\begin{itemize}
  \item \codeinline{feedback}: feedback del paciente con nivel de dolor (1--10), estado emocional (5 valores), URL de audio opcional y, en línea, la \codeinline{transcripcion} (Whisper) y el \codeinline{resumen\_ia} (Ollama). No existe una tabla separada de análisis de IA.

  \item \codeinline{puntos\_dolor}: zonas del mapa corporal marcadas junto al feedback (parte del cuerpo, lado, intensidad).

  \item \codeinline{escalas\_resultado}: mediciones clínicas registradas por el terapeuta (\codeinline{ROM}, \codeinline{STRENGTH}, \codeinline{FUNCTIONAL}, \codeinline{VAS}) con valor y unidad.
\end{itemize}

\subsection{Dominio de Agenda}

\begin{itemize}
  \item \codeinline{citas}: citas entre paciente y terapeuta con estado (\codeinline{SCHEDULED}, \codeinline{CONFIRMED}, \codeinline{CANCELLED}, \codeinline{COMPLETED}), sala y plan opcionales, y \codeinline{grupo\_recurrencia\_id} para series recurrentes.

  \item \codeinline{salas}: consultorios disponibles con nombre, ubicación, capacidad y equipamiento.

  \item \codeinline{lista\_espera}: inscripciones del paciente para avisarle cuando se libere un hueco con su terapeuta dentro de un rango de fechas.
\end{itemize}

\subsection{Dominio de Comunicación}

\begin{itemize}
  \item \codeinline{conversaciones} y \codeinline{mensajes}: hilo de mensajería en tiempo real entre un paciente y su terapeuta (única por par), con marca de leído.

  \item \codeinline{notificaciones}: centro de notificaciones in-app por usuario (tipo, título, cuerpo, enlace y estado de lectura).

  \item \codeinline{preferencias\_notificacion}: preferencias de correo, recordatorios y horario silencioso por usuario.
\end{itemize}

\subsection{Dominio de Plataforma}

\begin{itemize}
  \item \codeinline{auditoria}: bitácora de acciones sensibles (inicio de sesión, consulta de historial, subida de documentos) con usuario, entidad, IP y metadatos.

  \item \codeinline{tokens\_refresco} y \codeinline{tokens\_revocados}: soporte de la sesión JWT (refresh de 7 días y revocación al cerrar sesión).
\end{itemize}

\section{Diagrama ERD}

\begin{sidewaysfigure}
  \centering
  \includegraphics[width=\textheight,height=\textwidth,keepaspectratio]{FisioManager_ERD.png}
  \caption{Modelo de Base de Datos ERD de FisioManager: 30 tablas con claves primarias, claves foráneas y restricciones.}
  \label{fig:modelo-erd}
\end{sidewaysfigure}

\section{Convenciones aplicadas}

\begin{itemize}
  \item \textbf{Claves primarias (PK):} todas las tablas usan un \codeinline{id} de texto tipo \codeinline{cuid} generado por Prisma, excepto las extensiones de perfil (\codeinline{perfiles\_paciente}, \codeinline{perfiles\_terapeuta}) que usan \codeinline{usuario\_id} como PK y FK simultáneamente.

  \item \textbf{Claves foráneas (FK):} se definen restricciones de integridad referencial en todas las relaciones, con borrado en cascada donde corresponde (p. ej. perfiles, fases, mensajes).

  \item \textbf{Timestamps:} las tablas principales incluyen \codeinline{created\_at} y \codeinline{updated\_at} para auditoría temporal.

  \item \textbf{Valores acotados:} los campos enumerados (rol, estados, tipos) se almacenan como \codeinline{TEXT} y se validan en la capa de aplicación con esquemas Zod.

  \item \textbf{Unicidad:} \codeinline{email} en \codeinline{usuarios}, \codeinline{paciente\_id} en \codeinline{historial\_clinico}, el par (\codeinline{paciente\_id}, \codeinline{terapeuta\_id}) en \codeinline{conversaciones} y el par (\codeinline{terapeuta\_id}, \codeinline{dia\_semana}) en \codeinline{disponibilidad\_terapeuta} tienen restricción \codeinline{UNIQUE}.

  \item \textbf{Índices:} se añaden índices sobre las claves foráneas y los campos de filtrado frecuente (auditoría por usuario/entidad/fecha, notificaciones por usuario y estado de lectura).
\end{itemize}

\section{Código PlantUML del diagrama}

El código fuente se entrega como \codeinline{12\_modelo\_erd.puml} en el directorio \codeinline{puml/}. Se regenera con:

\begin{lstlisting}[language=bash]
plantuml -tpng 12_modelo_erd.puml
\end{lstlisting}

El código fuente completo se incluye en el Anexo~\ref{ch:anexos}.
